\documentclass[11pt,a4paper]{report}

% ===================== PACKAGES =====================
\usepackage[utf8]{inputenc}
\usepackage[T1]{fontenc}
\usepackage[french]{babel}
\usepackage{lmodern}
\usepackage[margin=2.5cm]{geometry}
\usepackage{xcolor}
\usepackage{tikz}
\usetikzlibrary{arrows.meta,positioning,shapes.multipart,fit,backgrounds,calc,shadows}
\usepackage{booktabs}
\usepackage{longtable}
\usepackage{tabularx}
\usepackage{array}
\usepackage{multirow}
\usepackage{colortbl}
\usepackage{enumitem}
\usepackage{titlesec}
\usepackage{fancyhdr}
\usepackage{rotating}
\usepackage{graphicx}
\usepackage{adjustbox}
\usepackage[skins,breakable]{tcolorbox}
\usepackage{hyperref}
\usepackage{parskip}

% ===================== CHAMPS À COMPLÉTER =====================
% Renseignez ces champs avant compilation (informations académiques).
\newcommand{\ecoleNom}{École Marocaine des Sciences de l'Ingénieur}
\newcommand{\ecoleSigle}{EMSI}
\newcommand{\villeEcole}{Casablanca --- Campus Oranger}
\newcommand{\filiereEtud}{Développement des Systèmes d'Information (DSI) --- 4\textsuperscript{e} année}
\newcommand{\anneeUniv}{2025 -- 2026}
\newcommand{\nomEtudiant}{Marouane Essaadouni}
\newcommand{\nomEncadrant}{Asmaa Roudane}
\newcommand{\typeProjet}{Stage de fin d'année}

% ===================== COULEURS =====================
\definecolor{km0blue}{RGB}{22,52,89}
\definecolor{km0orange}{RGB}{224,122,29}
\definecolor{km0green}{RGB}{45,134,89}
\definecolor{km0red}{RGB}{178,34,34}
\definecolor{km0gray}{RGB}{240,242,245}
\definecolor{km0graytext}{RGB}{90,98,110}

% ===================== TITRES =====================
\titleformat{\chapter}[hang]{\normalfont\Huge\bfseries\color{km0blue}}{\thechapter.}{20pt}{\Huge}
\titleformat{\section}{\normalfont\Large\bfseries\color{km0blue}}{\thesection}{10pt}{}
\titleformat{\subsection}{\normalfont\large\bfseries\color{km0blue}}{\thesubsection}{8pt}{}

% ===================== EN-TETES =====================
\pagestyle{fancy}
\fancyhf{}
\renewcommand{\headrulewidth}{0.4pt}
\fancyhead[L]{\small\textcolor{km0graytext}{KM0 --- Rapport de stage}}
\fancyhead[R]{\small\textcolor{km0graytext}{\leftmark}}
\fancyfoot[C]{\small\thepage}

% ===================== HYPERREF =====================
\hypersetup{
    colorlinks=true,
    linkcolor=km0blue,
    citecolor=km0blue,
    urlcolor=km0blue,
    pdftitle={KM0 - Rapport de stage},
    pdfauthor={Marouane Essaadouni}
}

% ===================== ENVIRONNEMENTS UTILES =====================
\newtcolorbox{regle}[2]{
  enhanced, breakable,
  colback=km0blue!4, colframe=km0blue,
  colbacktitle=km0blue, coltitle=white,
  fonttitle=\bfseries,
  title={Règle #1 --- #2},
  boxrule=0.6pt, arc=2pt, top=5pt, bottom=5pt, left=7pt, right=7pt,
  before skip=8pt, after skip=8pt
}

\newtcolorbox{avancement}[1]{
  enhanced, breakable,
  colback=km0orange!6, colframe=km0orange,
  fonttitle=\bfseries, coltitle=km0orange!40!black,
  title={État d'avancement --- #1},
  boxrule=0.6pt, arc=2pt, top=5pt, bottom=5pt, left=7pt, right=7pt,
  before skip=8pt, after skip=8pt
}

\newcommand{\etat}[3]{%
  % #1 = texte, #2 = couleur, #3 = position tikz optionnelle
  \tikz[baseline=-0.6ex]\node[rectangle,rounded corners=2pt,fill=#2!15,draw=#2,thick,inner sep=3pt]{\textcolor{#2!70!black}{\textbf{#1}}};%
}

\newcommand{\capture}[3]{%
  % #1 = fichier (sans extension), #2 = largeur, #3 = légende
  \begin{figure}[h!]
  \centering
  \fbox{\includegraphics[width=#2\linewidth]{screenshots/#1.png}}
  \caption{#3}
  \end{figure}%
}

\newcommand{\diagram}[3]{%
  % #1 = fichier (sans extension) dans diagrams/, #2 = largeur, #3 = légende
  % Image générée hors LaTeX (PlantUML) --- voir DIAGRAMMES.md à la racine du projet.
  \begin{figure}[h!]
  \centering
  \fbox{\includegraphics[width=#2\linewidth]{diagrams/#1.png}}
  \caption{#3}
  \end{figure}%
}

% ===================== DOCUMENT =====================
\begin{document}

% --------------------------------------------------------------
% PAGE DE GARDE
% --------------------------------------------------------------
\begin{titlepage}
\centering

% ---- Bandeau supérieur : logo ----
% Placez le fichier logo.png à la racine du projet Overleaf (même dossier que ce .tex).
\includegraphics[width=4.6cm]{logo.png}

\vspace{0.6cm}
{\large\bfseries\color{km0blue} \ecoleNom}\\[0.1cm]
{\normalsize\color{km0graytext} \ecoleSigle\ --- \villeEcole}\\[0.15cm]
{\small\color{km0graytext} Filière~: \filiereEtud}\\[0.1cm]
{\small\color{km0graytext}\bfseries \typeProjet}

\vspace{0.4cm}
\begin{tikzpicture}
\draw[km0blue, line width=1pt] (-7,0) -- (7,0);
\end{tikzpicture}

\vspace{2cm}
{\Huge\bfseries\color{km0blue} RAPPORT DE STAGE}\\[1.2cm]

{\huge\bfseries\color{km0blue} KM0}\\[0.25cm]
{\large\color{km0graytext} Conception et développement d'un système de gestion \\ de garage automobile}

\vspace{0.3cm}
\begin{tikzpicture}
\draw[km0orange, line width=1.4pt] (-1.8,0) -- (1.8,0);
\end{tikzpicture}

\vspace{0.4cm}
{\color{km0graytext} Version intermédiaire --- 7 août 2026}

\vfill

\begin{tikzpicture}
\draw[km0blue!40, line width=0.6pt] (-7,0) -- (7,0);
\end{tikzpicture}
\vspace{0.5cm}

\begin{tabularx}{0.85\linewidth}{@{}X X@{}}
\textbf{Réalisé par~:} & \textbf{Encadré par~:} \\[0.15cm]
\nomEtudiant & \nomEncadrant \\
\end{tabularx}

\vspace{1.2cm}
{\small\color{km0graytext} \villeEcole\ --- Année universitaire \anneeUniv}

\vspace{0.5cm}
\end{titlepage}

% --------------------------------------------------------------
\chapter*{Remerciements}
\addcontentsline{toc}{chapter}{Remerciements}

Je tiens à exprimer ma sincère gratitude à toutes les personnes qui ont contribué,
de près ou de loin, à la réalisation de ce projet de fin d'année.

Mes remerciements s'adressent en premier lieu à mon encadrante, \textbf{\nomEncadrant},
pour son suivi régulier, ses conseils méthodologiques et sa disponibilité tout au long
de ce travail.

Je remercie également l'équipe pédagogique de \ecoleSigle\ pour la qualité de la formation
reçue en \filiereEtud, qui m'a donné les bases nécessaires pour mener à bien la conception
et le développement d'une application de gestion complète, de l'analyse des besoins jusqu'à
l'implémentation technique.

Enfin, je remercie le garage ayant servi de terrain d'observation pour ce projet, dont le
fonctionnement quotidien --- et ses limites actuelles --- ont directement inspiré le
cahier des charges et les choix fonctionnels de l'application KM0.

\clearpage

% --------------------------------------------------------------
\chapter*{Résumé}
\addcontentsline{toc}{chapter}{Résumé}

KM0 est une application web de gestion destinée à un garage automobile marocain travaillant
à la fois pour des particuliers et des institutions étatiques. Le projet vise à remplacer un
suivi actuellement informel (papier, fichiers Excel, mémoire des employés) par un outil
centralisé, traçable et sécurisé, couvrant l'intégralité du cycle de vie d'un véhicule pris
en charge par le garage, depuis son entrée jusqu'à sa sortie.

Ce rapport présente la démarche complète du projet~: analyse du contexte métier et de la
problématique, spécification des besoins, conception de l'architecture technique et du
modèle de données, puis réalisation progressive des modules applicatifs. À la date de
rédaction, l'authentification par rôle, la gestion des clients, la gestion des véhicules
et la création de dossiers d'intervention (avec génération automatique de la fiche
d'intervention en PDF) sont opérationnelles. Les modules restants --- vue Kanban de
l'atelier, génération du Bon de Sortie sécurisé par QR code, écran de vérification à la
sortie, facturation et tableau de bord --- sont conçus et planifiés, et feront l'objet
d'un développement itératif documenté dans les prochaines versions de ce rapport.

\textbf{Mots-clés~:} gestion de garage automobile, Next.js, TypeScript, Prisma, PostgreSQL,
traçabilité, Bon de Sortie, QR code, contrôle d'accès par rôle.

\clearpage

\tableofcontents
\clearpage

% ================================================================
\chapter{Introduction générale}
% ================================================================

Les petites et moyennes structures artisanales, en particulier dans le secteur de la
réparation automobile, souffrent souvent d'un déficit d'outillage numérique adapté à leur
taille. Le suivi des véhicules pris en charge repose fréquemment sur des supports
hétérogènes --- carnets papier, fichiers tableur, mémoire des employés --- qui ne
garantissent ni traçabilité, ni sécurité, ni visibilité de gestion. Cette situation devient
particulièrement problématique lorsque l'activité mêle une clientèle de particuliers à une
clientèle institutionnelle (véhicules de fonction d'administrations publiques), pour
laquelle les exigences de formalisme et de justification sont plus élevées.

C'est dans ce contexte qu'a été initié le projet \textbf{KM0}, une application web de
gestion de garage automobile, développée dans le cadre du stage de fin d'année de la
filière Développement des Systèmes d'Information à l'\ecoleSigle. L'objectif du projet est
double~: d'une part, digitaliser et centraliser l'ensemble du processus métier du garage
(de l'entrée du véhicule à sa sortie) ; d'autre part, sécuriser formellement l'étape la
plus sensible du processus --- la sortie du véhicule --- à l'aide d'un document
vérifiable et à usage unique, le \emph{Bon de Sortie}.

Ce rapport documente l'ensemble de la démarche suivie~: analyse du contexte et de la
problématique, spécification fonctionnelle et non fonctionnelle, conception de
l'architecture logicielle et du modèle de données, puis réalisation technique des modules
applicatifs. Le développement étant organisé selon un planning progressif détaillé
(une tâche métier par jour ouvré, sur quatre semaines), ce rapport reflète l'état
d'avancement du projet à la date de rédaction et sera complété au fil des livraisons
successives~: les chapitres relatifs aux fonctionnalités non encore développées
présentent leur conception et seront enrichis, voire corrigés, une fois l'implémentation
réalisée.

Le document est organisé comme suit~: le chapitre~2 présente le projet et son organisme
d'accueil ; le chapitre~3 détaille l'analyse des besoins ; le chapitre~4 expose la
conception technique retenue ; le chapitre~5 décrit la réalisation des modules
actuellement opérationnels ; le chapitre~6 dresse un état d'avancement par rapport au
planning initial ; le chapitre~7 présente la démarche de test et de qualité ; enfin, le
chapitre~8 conclut ce rapport et ouvre sur les perspectives du projet.

% ================================================================
\chapter{Présentation du projet et de l'organisme d'accueil}
% ================================================================

\section{Présentation de l'organisme d'accueil}

\begin{table}[h!]
\centering
\rowcolors{2}{km0gray}{white}
\begin{tabularx}{\linewidth}{@{}l>{\raggedright\arraybackslash}X@{}}
\toprule
\rowcolor{km0blue}
\textcolor{white}{\textbf{Champ}} & \textcolor{white}{\textbf{Valeur}} \\
\midrule
Gérant & Marouane Essaadouni \\
Activité & Mécanicien réparateur \\
Adresse & 24, Allée des Tamaris, Aïn Sebaâ, Casablanca, Maroc \\
Registre de Commerce (RC) & 397421 \\
Tribunal de commerce & Casablanca \\
ICE & 002027113000005 \\
Capital social & 100 000 DHS \\
\bottomrule
\end{tabularx}
\caption{Fiche signalétique de l'entreprise d'accueil}
\end{table}

\section{Contexte métier}

\textbf{KM0} est une application web de gestion destinée à un garage automobile marocain
travaillant à la fois pour des \textbf{particuliers} et pour des \textbf{institutions
étatiques}. Elle a pour objectif de remplacer un suivi actuellement informel par un outil
centralisé, traçable et sécurisé, couvrant l'intégralité du cycle de vie d'un véhicule
pris en charge par le garage~: de son entrée jusqu'à sa sortie.

Le garage traite quatre grandes familles d'interventions~:

\begin{table}[h!]
\centering
\rowcolors{2}{km0gray}{white}
\begin{tabularx}{\linewidth}{@{}l>{\raggedright\arraybackslash}X@{}}
\toprule
\rowcolor{km0blue}
\textcolor{white}{\textbf{Type d'intervention}} & \textcolor{white}{\textbf{Description}} \\
\midrule
\textbf{Diagnostic} & Identification d'une panne ou d'un dysfonctionnement. \\
\textbf{Réparation} & Intervention mécanique~: moteur, transmission, etc. \\
\textbf{Entretien} & Vidange, changement de filtres, révision périodique. \\
\textbf{Tôlerie / Carrosserie} & Véhicules accidentés, redressage, peinture. \\
\bottomrule
\end{tabularx}
\caption{Types d'interventions couvertes par le système}
\end{table}

Un même dossier peut cumuler plusieurs types d'intervention (par exemple un diagnostic
suivi d'une réparation) --- ce cas est directement pris en charge dans le modèle de
données et dans le formulaire de création de dossier (voir chapitre~5).

\section{Problématique}

Le fonctionnement actuel du garage repose sur des outils non structurés, ce qui engendre
plusieurs difficultés opérationnelles~:

\begin{table}[h!]
\centering
\rowcolors{2}{km0gray}{white}
\begin{tabularx}{\linewidth}{@{}>{\raggedright\arraybackslash}X>{\raggedright\arraybackslash}X@{}}
\toprule
\rowcolor{km0blue}
\textcolor{white}{\textbf{Problème constaté}} & \textcolor{white}{\textbf{Conséquence}} \\
\midrule
Suivi papier / Excel / mémoire des employés & Aucune traçabilité, perte d'information \\
Pas de contrôle formalisé à la sortie du véhicule & Véhicules sortis sans paiement ou sans autorisation \\
Historique client dispersé & Impossible de savoir ce qui a été fait sur un véhicule \\
Aucune visibilité pour le gérant & Pas de vue sur la charge de travail, le chiffre d'affaires, les délais \\
\bottomrule
\end{tabularx}
\caption{Problèmes actuels et conséquences}
\end{table}

\section{Objectifs du projet}

Le projet KM0 poursuit quatre objectifs principaux~:

\begin{enumerate}[label=\arabic*., leftmargin=1.2cm]
\item \textbf{Centraliser} la gestion des dossiers véhicules, des clients et des interventions.
\item \textbf{Sécuriser la sortie} des véhicules via un \emph{Bon de Sortie} (BS) vérifiable,
      empêchant toute sortie non autorisée ou non payée.
\item \textbf{Offrir à toute l'équipe une meilleure traçabilité et organisation} du travail
      quotidien~: chaque intervenant (réceptionniste, technicien, admin) sait à tout moment
      ce qui doit être fait et ce qui a déjà été fait sur un dossier.
\item \textbf{Donner au gérant une vision business exploitable}~: charge de l'atelier,
      chiffre d'affaires, délais de traitement.
\end{enumerate}

% ================================================================
\chapter{Analyse et spécification des besoins}
% ================================================================

\section{Acteurs du système}

Le contrôle d'accès repose sur quatre rôles distincts, appliqués à deux niveaux~: via le
middleware Next.js pour le routage, et systématiquement à l'intérieur de chaque Server
Action (défense en profondeur).

\begin{table}[h!]
\centering
\rowcolors{2}{km0gray}{white}
\begin{tabularx}{\linewidth}{@{}l>{\raggedright\arraybackslash}X@{}}
\toprule
\rowcolor{km0blue}
\textcolor{white}{\textbf{Rôle}} & \textcolor{white}{\textbf{Permissions}} \\
\midrule
\textbf{RECEPTIONNISTE} & Créer et éditer clients, véhicules, dossiers. Voir tous les dossiers.
Encaisser les paiements. \textbf{Ne peut pas} générer de Bon de Sortie. \\
\textbf{TECHNICIEN} & Voir les dossiers qui lui sont assignés. Changer le statut
(\texttt{EN\_DIAGNOSTIC} $\rightarrow$ \texttt{EN\_COURS} $\rightarrow$ \texttt{PRET}).
Ajouter des observations et des pièces utilisées. \\
\textbf{ADMIN} & Accès total. Génération du Bon de Sortie. Dérogations de paiement.
Consultation des statistiques. Gestion des utilisateurs. Régénération d'un BS perdu. \\
\textbf{VIGILE} & Accès \textbf{exclusif} à l'écran \texttt{/verification}. Redirection forcée
depuis toute autre route. Aucune donnée financière visible. \\
\bottomrule
\end{tabularx}
\caption{Matrice des rôles et permissions}
\end{table}

\section{Besoins fonctionnels}

Les besoins fonctionnels couvrent l'ensemble du cycle de vie d'un dossier~:

\begin{figure}[h!]
\centering
\begin{adjustbox}{max width=\linewidth}
\begin{tikzpicture}[
  step/.style={rectangle, rounded corners=3pt, draw=km0blue, thick, fill=km0blue!7,
               minimum width=2.9cm, minimum height=1.1cm, align=center, font=\scriptsize},
  final/.style={rectangle, rounded corners=3pt, draw=km0green, thick, fill=km0green!12,
               minimum width=2.9cm, minimum height=1.1cm, align=center, font=\scriptsize},
  arr/.style={-Stealth, thick, km0graytext}
]
\node[step] (arrivee) {Arrivée du véhicule};
\node[step, right=1cm of arrivee] (dossier) {Création du dossier\\(réceptionniste)};
\node[step, right=1cm of dossier] (fiche) {Fiche d'intervention\\générée};
\node[step, right=1cm of fiche] (charge) {Prise en charge\\(technicien)};
\node[step, right=1cm of charge] (avance) {Avancement\\(Kanban)};
\node[final, right=1cm of avance] (sortie) {BS $\rightarrow$ Sortie};

\draw[arr] (arrivee) -- (dossier);
\draw[arr] (dossier) -- (fiche);
\draw[arr] (fiche) -- (charge);
\draw[arr] (charge) -- (avance);
\draw[arr] (avance) -- (sortie);
\end{tikzpicture}
\end{adjustbox}
\caption{Logique de bout en bout du traitement d'un véhicule}
\end{figure}

\begin{itemize}[leftmargin=1cm]
\item Authentification et gestion des rôles.
\item Gestion des clients (particulier, institution étatique, entreprise) et de leurs véhicules.
\item Création de dossiers d'intervention, avec sélection multiple des types d'intervention.
\item Génération automatique d'une fiche d'intervention imprimable dès la création du dossier.
\item Suivi de l'avancement des dossiers via une vue Kanban de l'atelier.
\item Génération d'un Bon de Sortie sécurisé (PDF + QR code), bloqué tant que le dossier
      n'est pas au statut \texttt{PRET} et que le paiement n'est pas soldé (sauf dérogation).
\item Vérification de la sortie par scan du Bon de Sortie sur un écran dédié.
\item Facturation et suivi des paiements.
\item Tableau de bord de synthèse pour le gérant.
\item Historique complet par véhicule et par dossier.
\end{itemize}

\subsection{Diagramme de cas d'utilisation}

% Image générée hors LaTeX (PlantUML) --- voir DIAGRAMMES.md à la racine du projet.
% Fichier attendu : diagrams/diagramme_cas_utilisation.png
Le diagramme suivant formalise, pour chaque acteur du système, les cas d'utilisation
auxquels il a accès, en cohérence avec la matrice des rôles et permissions présentée
en section~3.1.

\diagram{diagramme_cas_utilisation}{0.95}{Diagramme de cas d'utilisation du système KM0}

\section{Besoins non fonctionnels}

\begin{itemize}[leftmargin=1cm]
\item \textbf{Sécurité} --- authentification par session (cookies \texttt{httpOnly}, mots de
      passe hachés avec \texttt{bcrypt}), contrôle d'accès vérifié à double niveau
      (middleware et Server Actions).
\item \textbf{Traçabilité} --- tout changement de statut d'un dossier est enregistré avec
      l'utilisateur responsable, de manière non contournable depuis l'interface.
\item \textbf{Ergonomie} --- l'écran de vérification de sortie, utilisé par un agent de
      sécurité non technique sur tablette, doit rester lisible à distance et ne comporter
      aucun jargon technique.
\item \textbf{Autonomie} --- aucune dépendance à un service cloud tiers pour
      l'authentification, afin de permettre une évolution future vers un déploiement en
      réseau local (hors périmètre du présent projet).
\item \textbf{Fiabilité du typage} --- développement en TypeScript strict, sans type
      \texttt{any}, avec vérification systématique (\texttt{tsc --noEmit}) avant chaque
      livraison.
\item \textbf{Intégrité transactionnelle} --- toute opération sensible à la concurrence
      (numérotation séquentielle, invalidation d'un Bon de Sortie) est exécutée dans une
      transaction de base de données garantissant l'absence de collision.
\end{itemize}

\section{Périmètre du projet}

\subsection{Fonctionnalités couvertes (MVP)}

\begin{itemize}[leftmargin=1cm]
\item Authentification et gestion des rôles
\item Gestion des clients et des véhicules
\item Gestion des dossiers et vue Kanban de l'atelier
\item Génération automatique d'une fiche d'intervention à la création du dossier
\item Génération du Bon de Sortie (PDF + QR code)
\item Écran de vérification de sortie (scan QR)
\item Facturation
\item Tableau de bord de synthèse
\item Historique complet par véhicule
\end{itemize}

\subsection{Hors périmètre (V2, non développé sauf demande explicite)}

Afin de conserver un périmètre réaliste pour la durée du stage, les fonctionnalités
suivantes sont volontairement exclues de la version actuelle, sans pour autant être
ignorées dans les choix d'architecture (les services métier restent découplés pour
faciliter une migration future)~:

\begin{itemize}[leftmargin=1cm]
\item Gestion de stock détaillée des pièces détachées
\item Planning technicien avancé (affectation par créneaux, charge prévisionnelle)
\item Notifications SMS / WhatsApp réelles vers les clients
\item Mode de fonctionnement en réseau local sans connexion Internet
\item Système de licence multi-postes
\item Gestion multi-garage
\item Application mobile native
\end{itemize}

% ================================================================
\chapter{Conception}
% ================================================================

\section{Choix technologiques}

\begin{table}[h!]
\centering
\rowcolors{2}{km0gray}{white}
\begin{tabularx}{\linewidth}{@{}l>{\raggedright\arraybackslash}X@{}}
\toprule
\rowcolor{km0blue}
\textcolor{white}{\textbf{Composant}} & \textcolor{white}{\textbf{Technologie}} \\
\midrule
Framework applicatif & Next.js 15 (App Router) + TypeScript (mode strict) \\
Base de données & PostgreSQL 16 \\
ORM & Prisma 6 \\
Interface utilisateur & Tailwind CSS + shadcn/ui (primitives Base UI) \\
Authentification & Session personnalisée (cookies \texttt{httpOnly} + \texttt{bcrypt}) \\
Génération PDF & \texttt{@react-pdf/renderer} \\
QR Code & \texttt{qrcode} (génération) + \texttt{html5-qrcode} (lecture navigateur) \\
Validation des données & Zod \\
Gestion des dates & \texttt{date-fns} (locale \texttt{fr}) \\
Tests & Vitest \\
\bottomrule
\end{tabularx}
\caption{Stack technique retenue}
\end{table}

Ce choix technologique répond à plusieurs critères~: un framework full-stack unique
(Next.js) permettant de mutualiser rendu serveur, routage et logique applicative ; un ORM
typé (Prisma) garantissant la cohérence entre le schéma de base de données et les types
TypeScript utilisés dans le code applicatif ; une validation systématique des entrées
utilisateur (Zod) à la frontière de chaque action serveur ; et l'absence de dépendance à
un fournisseur d'authentification cloud, condition posée dès le cahier des charges pour
garder l'application autonome.

\section{Contraintes techniques}

\begin{itemize}[leftmargin=1cm]
\item \textbf{Aucune dépendance cloud tierce} pour l'authentification (pas de Clerk,
      Auth0, Supabase Auth)~: l'application doit rester autonome.
\item Les \textbf{Server Actions} sont privilégiées par rapport aux routes API. Les seules
      exceptions sont l'endpoint de vérification du Bon de Sortie, qui doit rester
      accessible via une route API classique pour être appelable par le scanner, et les
      endpoints de génération de documents PDF (fiche d'intervention, Bon de Sortie), qui
      nécessitent une route API pour renvoyer un flux binaire téléchargeable.
\item Aucun type \texttt{any} n'est toléré en TypeScript ; la commande
      \texttt{tsc --noEmit} doit systématiquement passer sans erreur avant chaque
      livraison.
\end{itemize}

\section{Architecture applicative}

\begin{figure}[h!]
\centering
\begin{adjustbox}{max width=\linewidth}
\begin{tikzpicture}[
  box/.style={rectangle, rounded corners=3pt, draw=km0blue, thick, fill=km0blue!6,
              minimum width=3.4cm, minimum height=1.1cm, align=center, font=\small},
  boxalt/.style={rectangle, rounded corners=3pt, draw=km0orange, thick, fill=km0orange!8,
              minimum width=3.4cm, minimum height=1.1cm, align=center, font=\small},
  db/.style={cylinder, shape border rotate=90, aspect=0.25, draw=km0blue, thick,
             fill=km0blue!10, minimum width=2.6cm, minimum height=1.4cm, align=center, font=\small},
  arr/.style={-Stealth, thick, km0graytext}
]
\node[box] (client) {Navigateur\\(réceptionniste, technicien, admin)};
\node[boxalt, right=2.4cm of client] (vigile) {Tablette guérite\\(vigile)};

\node[box, below=1.3cm of $(client)!0.5!(vigile)$, minimum width=6.5cm] (next) {Next.js 15 --- App Router\\Server Actions};
\node[boxalt, below=0.55cm of next, minimum width=6.5cm] (api) {Routes API (\texttt{verification}, \texttt{fiche PDF})};

\node[box, below=0.9cm of api, xshift=-2.6cm] (prisma) {Prisma ORM};
\node[box, below=0.9cm of api, xshift=2.6cm] (services) {Services métier\\(\texttt{lib/services})};

\node[db, below=1.1cm of $(prisma)!0.5!(services)$] (pg) {PostgreSQL 16};

\node[boxalt, right=2.8cm of services] (pdf) {PDF + QR\\(fiche, BS)};

\draw[arr] (client) -- (next);
\draw[arr] (vigile) -- (api);
\draw[arr] (next) -- (api);
\draw[arr] (next) -- (services);
\draw[arr] (api) -- (services);
\draw[arr] (services) -- (prisma);
\draw[arr] (prisma) -- (pg);
\draw[arr] (services) -- (pdf);
\end{tikzpicture}
\end{adjustbox}
\caption{Architecture applicative simplifiée}
\end{figure}

L'architecture repose sur une séparation stricte des responsabilités~: les composants
React (pages et formulaires) ne contiennent aucune logique métier ; celle-ci est
entièrement déléguée à une couche de services (\texttt{lib/services/}), elle-même seule
autorisée à dialoguer avec la base de données via Prisma. Les Server Actions jouent le
rôle de contrôleur~: elles valident les entrées (Zod), vérifient les droits de
l'utilisateur (\texttt{requireRole}), puis délèguent au service métier correspondant.

\section{Modélisation des données}

\begin{sidewaysfigure}[p]
\centering
\resizebox{\textheight}{!}{%
\begin{tikzpicture}[
  entity/.style={rectangle, draw=km0blue, thick, fill=white, rounded corners=2pt,
                 align=left, font=\footnotesize, drop shadow},
  title/.style={font=\bfseries\small\color{white}, fill=km0blue},
  rel/.style={-Stealth, thick, km0graytext, font=\scriptsize},
  every node/.style={inner sep=4pt}
]

% ---- User ----
\node[entity] (user) at (0,7) {
  \begin{tabular}{@{}l@{}}
  \multicolumn{1}{c}{\cellcolor{km0blue}\textcolor{white}{\textbf{User}}}\\
  id, nom, email\\ passwordHash, role\\ actif
  \end{tabular}
};

% ---- Client ----
\node[entity] (client) at (11,9) {
  \begin{tabular}{@{}l@{}}
  \multicolumn{1}{c}{\cellcolor{km0blue}\textcolor{white}{\textbf{Client}}}\\
  id, nom, telephone\\ email, cin, type\\ raisonSociale, ice
  \end{tabular}
};

% ---- Vehicule ----
\node[entity] (vehicule) at (11,6) {
  \begin{tabular}{@{}l@{}}
  \multicolumn{1}{c}{\cellcolor{km0blue}\textcolor{white}{\textbf{Vehicule}}}\\
  id, immatriculation\\ marque, modele, annee\\ couleur, kilometrage
  \end{tabular}
};

% ---- Dossier (hub central) ----
\node[entity, fill=km0blue!8] (dossier) at (11,2.5) {
  \begin{tabular}{@{}l@{}}
  \multicolumn{1}{c}{\cellcolor{km0blue}\textcolor{white}{\textbf{Dossier}}}\\
  id, numero\\ typesIntervention[]\\ statut, motifDeclare\\
  dateEntree, datePrevueSortie\\ dateSortieReelle\\
  montantTotal, statutPaiement
  \end{tabular}
};

% ---- BonDeSortie ----
\node[entity] (bs) at (0,3.6) {
  \begin{tabular}{@{}l@{}}
  \multicolumn{1}{c}{\cellcolor{km0blue}\textcolor{white}{\textbf{BonDeSortie}}}\\
  id, numero, qrToken\\ statut\\ dateGeneration\\ dateUtilisation\\
  derogationPaiement\\ derogationMotif
  \end{tabular}
};

% ---- Facture ----
\node[entity] (facture) at (20,4.5) {
  \begin{tabular}{@{}l@{}}
  \multicolumn{1}{c}{\cellcolor{km0blue}\textcolor{white}{\textbf{Facture}}}\\
  id, numero\\ montantHT, tauxTVA\\ montantTTC\\
  statutPaiement\\ modePaiement, datePaiement
  \end{tabular}
};

% ---- PieceUtilisee ----
\node[entity] (piece) at (20,1.6) {
  \begin{tabular}{@{}l@{}}
  \multicolumn{1}{c}{\cellcolor{km0blue}\textcolor{white}{\textbf{PieceUtilisee}}}\\
  id, nomPiece, reference\\ quantite, prixUnitaire
  \end{tabular}
};

% ---- HistoriqueStatut ----
\node[entity] (hist) at (11,-0.8) {
  \begin{tabular}{@{}l@{}}
  \multicolumn{1}{c}{\cellcolor{km0blue}\textcolor{white}{\textbf{HistoriqueStatut}}}\\
  id, statutAvant, statutApres\\ commentaire, createdAt
  \end{tabular}
};

% ---- Relations ----
\draw[rel] (client) -- (vehicule) node[midway, above, sloped]{1 -- N};
\draw[rel] (vehicule) -- (dossier) node[midway, right]{1 -- N};
\draw[rel] (dossier) -- (bs) node[midway, above, sloped]{1 -- N};
\draw[rel] (dossier) -- (facture) node[midway, above, sloped]{1 -- N};
\draw[rel] (dossier) -- (piece) node[midway, below, sloped]{1 -- N};
\draw[rel] (dossier) -- (hist) node[midway, right]{1 -- N};
\draw[rel] (user) -- (dossier) node[midway, above, sloped, xshift=-0.4cm]{1 -- N (technicien, optionnel)};
\draw[rel] (user) -- (hist) node[midway, below, sloped]{1 -- N (auteur)};
\draw[rel] (user) -- (bs) node[midway, left]{1 -- N (dérogation, optionnel)};

\end{tikzpicture}%
}
\caption{Diagramme entité-relation du modèle KM0}
\end{sidewaysfigure}

\begin{table}[h!]
\centering
\rowcolors{2}{km0gray}{white}
\begin{tabularx}{\linewidth}{@{}l>{\raggedright\arraybackslash}X@{}}
\toprule
\rowcolor{km0blue}
\textcolor{white}{\textbf{Entité}} & \textcolor{white}{\textbf{Rôle}} \\
\midrule
\textbf{User} & Compte utilisateur interne (réceptionniste, technicien, admin, vigile). \\
\textbf{Client} & Particulier, entreprise ou institution étatique propriétaire de véhicules. \\
\textbf{Vehicule} & Véhicule rattaché à un client, identifié par son immatriculation. \\
\textbf{Dossier} & Intervention en cours ou passée sur un véhicule ; entité centrale du système. \\
\textbf{BonDeSortie} & Autorisation de sortie du véhicule, matérialisée par un QR code unique. \\
\textbf{Facture} & Document de facturation rattaché à un dossier. \\
\textbf{PieceUtilisee} & Pièce détachée consommée dans le cadre d'un dossier. \\
\textbf{HistoriqueStatut} & Trace d'audit de chaque changement de statut d'un dossier. \\
\bottomrule
\end{tabularx}
\caption{Rôle de chaque entité du modèle}
\end{table}

\subsection{Diagramme de classes UML}

% Image générée hors LaTeX (PlantUML) --- voir DIAGRAMMES.md à la racine du projet.
% Fichier attendu : diagrams/diagramme_classes.png
Le diagramme entité-relation ci-dessus donne une vue relationnelle du modèle. La figure
suivante en propose la formalisation en \textbf{diagramme de classes UML}, avec les
attributs typés, les énumérations associées et les cardinalités des associations, telle
qu'elle a directement guidé l'écriture du schéma Prisma (chapitre~4).

\diagram{diagramme_classes}{0.95}{Diagramme de classes UML du modèle KM0}

\section{Règles de gestion critiques}

Les règles suivantes sont non négociables~: elles conditionnent l'intégrité et la
sécurité du processus de sortie des véhicules. Certaines sont déjà appliquées dans les
modules réalisés (règle~6, règle~7) ; les autres relèvent des modules en cours de
conception (Bon de Sortie, écran de vérification).

\begin{regle}{1}{Statut \texttt{SORTI} automatique uniquement}
Le statut \texttt{SORTI} n'est posé \textbf{que} via le scan du Bon de Sortie sur l'écran
\texttt{/verification}. Aucune interface ne permet de le poser manuellement.
\end{regle}

\begin{regle}{2}{Génération du BS conditionnée}
Le bouton de génération n'est visible que si le dossier est au statut \texttt{PRET}.
Si le statut de paiement n'est pas \texttt{SOLDE}, la génération est bloquée, sauf
dérogation d'un Admin (motif obligatoire, 10 caractères minimum, tracé dans le système).
\end{regle}

\begin{regle}{3}{Usage unique du BS}
Au moment du scan, le passage au statut \texttt{UTILISE} est effectué de manière atomique
(transaction Prisma). Un même QR code ne peut jamais autoriser deux sorties, même en cas
de double scan simultané.
\end{regle}

\begin{regle}{4}{Régénération d'un BS perdu}
En cas de perte, seul un Admin peut générer un nouveau Bon de Sortie ; l'ancien passe au
statut \texttt{ANNULE} dans la même transaction.
\end{regle}

\begin{regle}{5}{Multi-intervention, BS unique}
Un dossier peut cumuler plusieurs types d'intervention (\texttt{typesIntervention}), mais
ne génère jamais qu'un seul Bon de Sortie.
\end{regle}

\begin{regle}{6}{Audit trail obligatoire}
Tout changement de statut est écrit dans \texttt{HistoriqueStatut} avec l'utilisateur
responsable. Ce mécanisme est implémenté au niveau du service métier et n'est pas
contournable depuis l'interface. \emph{Règle en application depuis la création des
dossiers (voir chapitre~5).}
\end{regle}

\begin{regle}{7}{Numérotation séquentielle}
Les identifiants \texttt{DOS-2026-XXXX}, \texttt{BS-2026-XXXX} et \texttt{FACT-2026-XXXX}
sont générés séquentiellement par année, sans trou, via une transaction dédiée afin
d'éviter toute collision. \emph{Règle en application pour les dossiers (voir §5.4).}
\end{regle}

\section{Machine à états du dossier}

\begin{figure}[h!]
\centering
\begin{adjustbox}{max width=\linewidth}
\begin{tikzpicture}[
  st/.style={rectangle, rounded corners=3pt, draw=km0blue, thick, fill=km0blue!7,
             minimum width=2.6cm, minimum height=0.9cm, align=center, font=\scriptsize},
  final/.style={rectangle, rounded corners=3pt, draw=km0green, thick, fill=km0green!12,
             minimum width=2.6cm, minimum height=0.9cm, align=center, font=\scriptsize},
  cancel/.style={rectangle, rounded corners=3pt, draw=km0red, thick, fill=km0red!10,
             minimum width=2.6cm, minimum height=0.9cm, align=center, font=\scriptsize},
  tr/.style={-Stealth, thick, km0graytext, font=\tiny},
  trc/.style={-Stealth, thick, km0red!70, dashed, font=\tiny}
]
\node[st] (entre) {ENTRE};
\node[st, right=1.1cm of entre] (diag) {EN\_DIAGNOSTIC};
\node[st, right=1.1cm of diag] (devis) {DEVIS\_ENVOYE};
\node[st, above right=0.7cm and 1.3cm of devis] (accepte) {DEVIS\_ACCEPTE};
\node[st, below right=0.7cm and 1.3cm of devis] (refuse) {DEVIS\_REFUSE};
\node[st, right=1.6cm of accepte] (cours) {EN\_COURS};
\node[st, right=1.1cm of cours] (pret) {PRET};
\node[final, right=1.4cm of pret] (sorti) {SORTI};
\node[cancel, below=1.6cm of devis] (annule) {ANNULE};

\draw[tr] (entre) -- (diag);
\draw[tr] (diag) -- (devis);
\draw[tr] (devis) -- (accepte);
\draw[tr] (devis) -- (refuse);
\draw[tr] (accepte) -- (cours);
\draw[tr] (cours) -- (pret);
\draw[tr] (pret) -- node[above, font=\tiny\bfseries, text=km0green]{scan BS} (sorti);

\draw[trc] (diag) -- (annule);
\draw[trc] (devis) -- (annule);
\draw[trc] (refuse) -- (annule);
\draw[trc] (cours) -- (annule);
\end{tikzpicture}
\end{adjustbox}
\caption{Transitions de statut autorisées pour un dossier (traits pleins = flux normal,
traits pointillés rouges = annulation)}
\end{figure}

\begin{avancement}{machine à états}
À la date de rédaction, seule la transition initiale (création $\rightarrow$ \texttt{ENTRE})
est implémentée, avec écriture dans \texttt{HistoriqueStatut}. Le service de transition
générique (validation des transitions autorisées, assignation technicien) est planifié
pour la deuxième semaine de développement (voir chapitre~6).
\end{avancement}

\section{Écran de vérification --- spécification}

\begin{avancement}{écran de vérification}
Module non encore développé (planifié semaine 2). Sa spécification est présentée ici car
elle conditionne des choix de conception déjà retenus dans le modèle de données
(\texttt{BonDeSortie.qrToken}, \texttt{statut}).
\end{avancement}

L'écran \texttt{/verification} est conçu comme l'écran le plus critique du système. Il
sera utilisé par un agent de sécurité non technique, sur tablette, dans une guérite. Son
ergonomie devra donc privilégier la lisibilité et la simplicité absolue~: champ de saisie
en auto-focus combiné à un scanner QR par caméra, typographie très large lisible à
2~mètres, zéro jargon technique, retour automatique à l'état de scan après 8 secondes.

\begin{table}[h!]
\centering
\rowcolors{2}{km0gray}{white}
\begin{tabularx}{\linewidth}{@{}l>{\raggedright\arraybackslash}X@{}}
\toprule
\rowcolor{km0blue}
\textcolor{white}{\textbf{État}} & \textcolor{white}{\textbf{Affichage}} \\
\midrule
\etat{VERT}{km0green} & « SORTIE AUTORISÉE » --- immatriculation en très grand, marque/modèle, client. \\
\etat{ORANGE}{km0orange} & « SORTIE AUTORISÉE --- DÉROGATION » --- « Non payé, autorisé par [Admin] ». \\
\etat{ROUGE}{km0red} & « SORTIE REFUSÉE » avec raison précise~: facture non réglée, bon déjà
utilisé (avec date), bon inconnu, ou véhicule non prêt. \\
\bottomrule
\end{tabularx}
\caption{États d'affichage prévus de l'écran de vérification}
\end{table}

\subsection{Diagramme de séquence --- génération et vérification du BS}

\begin{figure}[h!]
\centering
\begin{adjustbox}{max width=\linewidth}
\begin{tikzpicture}[
  actor/.style={rectangle, draw=km0blue, thick, fill=km0blue!8, rounded corners=2pt,
                minimum width=2.6cm, minimum height=0.8cm, align=center, font=\scriptsize\bfseries},
  lifeline/.style={dashed, km0graytext},
  msg/.style={-Stealth, thick, km0blue, font=\scriptsize},
  msgret/.style={-Stealth, thick, km0graytext, dashed, font=\scriptsize},
  note/.style={rectangle, draw=km0orange, fill=km0orange!8, rounded corners=2pt, font=\scriptsize, align=left, inner sep=4pt}
]
% Acteurs
\node[actor] (admin) at (0,0) {Admin};
\node[actor] (app) at (4.2,0) {Application\\(bs.service)};
\node[actor] (bdd) at (8.4,0) {Base de\\données};
\node[actor] (vigile) at (12.6,0) {Vigile};

% Lifelines
\draw[lifeline] (admin.south) -- ++(0,-11.5);
\draw[lifeline] (app.south) -- ++(0,-11.5);
\draw[lifeline] (bdd.south) -- ++(0,-11.5);
\draw[lifeline] (vigile.south) -- ++(0,-11.5);

% Phase 1 : génération
\draw[msg] ($(admin.south)+(0,-1)$) -- node[above]{demande génération BS (dossier PRET)} ($(app.south)+(0,-1)$);
\draw[msg] ($(app.south)+(0,-1.9)$) -- node[above]{vérifie statutPaiement} ($(bdd.south)+(0,-1.9)$);
\draw[msgret] ($(bdd.south)+(0,-2.7)$) -- node[above]{SOLDE (ou dérogation)} ($(app.south)+(0,-2.7)$);
\draw[msg] ($(app.south)+(0,-3.6)$) -- node[above]{crée BonDeSortie (qrToken, VALIDE)} ($(bdd.south)+(0,-3.6)$);
\draw[msgret] ($(app.south)+(0,-4.6)$) -- node[above]{PDF A5 + QR code} ($(admin.south)+(0,-4.6)$);

\node[note, align=center] at (4.2,-5.4) {\textit{-- temps écoulé : véhicule prêt à récupérer --}};

% Phase 2 : vérification
\draw[msg] ($(vigile.south)+(0,-6.4)$) -- node[above]{scan QR (qrToken)} ($(app.south)+(0,-6.4)$);

\node[note] at (7,-7.6) {\textbf{Transaction atomique}\\BS existe ? statut VALIDE ?\\dossier PRET ? paiement soldé ?};

\draw[msg] ($(app.south)+(0,-9)$) -- node[above]{lecture + vérifications} ($(bdd.south)+(0,-9)$);
\draw[msg] ($(app.south)+(0,-9.8)$) -- node[above]{BS $\rightarrow$ UTILISE, Dossier $\rightarrow$ SORTI} ($(bdd.south)+(0,-9.8)$);
\draw[msgret] ($(app.south)+(0,-10.8)$) -- node[above]{VERT / ORANGE / ROUGE} ($(vigile.south)+(0,-10.8)$);

\end{tikzpicture}
\end{adjustbox}
\caption{Séquence complète : génération du BS puis vérification à la sortie (conception, développement prévu semaine 2)}
\end{figure}

% ================================================================
\chapter{Réalisation}
% ================================================================

Ce chapitre décrit les modules effectivement implémentés et opérationnels à la date de
rédaction de ce rapport. Chaque section présente l'objectif fonctionnel du module, les
choix techniques retenus et une illustration de l'interface obtenue.

\section{Organisation du code}

\begin{figure}[h!]
\centering
\begin{adjustbox}{max width=\linewidth}
\begin{tikzpicture}
\node[draw=km0blue, fill=km0gray, rounded corners=3pt, inner sep=10pt, align=left, font=\ttfamily\footnotesize]{%
\begin{tabular}{@{}l@{}}
app/ \\
\quad (auth)/login/ \\
\quad (app)/dashboard/ \\
\quad (app)/atelier/ \hfill\textcolor{km0graytext}{\% vue Kanban} \\
\quad (app)/dossiers/, dossiers/new/, dossiers/[id]/ \\
\quad (app)/clients/ \\
\quad (app)/vehicules/ \\
\quad (app)/factures/ \\
\quad (app)/admin/ \\
\quad verification/ \hfill\textcolor{km0graytext}{\% hors layout app, accès vigile} \\
\quad api/verification/route.ts \hfill\textcolor{km0graytext}{\% endpoint scan (POST)} \\
\quad api/dossiers/[id]/fiche/route.ts \hfill\textcolor{km0graytext}{\% PDF fiche d'intervention} \\
lib/ \\
\quad auth/ \hfill\textcolor{km0graytext}{\% session, hashing, guards} \\
\quad services/ \hfill\textcolor{km0graytext}{\% logique métier} \\
\quad validations/ \hfill\textcolor{km0graytext}{\% schémas Zod} \\
\quad utils/ \\
components/ \\
\quad ui/ \hfill\textcolor{km0graytext}{\% shadcn} \\
\quad features/ \hfill\textcolor{km0graytext}{\% composants métier} \\
prisma/ \\
\end{tabular}
};
\end{tikzpicture}
\end{adjustbox}
\caption{Arborescence de dossiers du projet}
\end{figure}

\begin{itemize}[leftmargin=1cm]
\item La logique métier réside exclusivement dans \texttt{lib/services/}, jamais dans les
      composants React.
\item Toute Server Action valide ses entrées avec Zod avant traitement.
\item Toute Server Action vérifie le rôle de l'utilisateur en première ligne.
\item Nommage en français pour le domaine métier (\texttt{dossier}, \texttt{vehicule}),
      en anglais pour la technique (\texttt{getUser}, \texttt{formatDate}).
\item Commits au format conventionnel~: \texttt{feat:}, \texttt{fix:}, \texttt{chore:},
      \texttt{test:}, \texttt{docs:}, un commit atomique par tâche.
\end{itemize}

\section{Authentification et gestion des rôles}

Le module d'authentification repose sur un mécanisme de session personnalisé, sans
dépendance externe. À la connexion, le mot de passe fourni est comparé à l'empreinte
stockée en base via \texttt{bcrypt} ; en cas de succès, un jeton de session est signé avec
l'algorithme HMAC-SHA256 (Web Crypto API) et déposé dans un cookie \texttt{httpOnly},
\texttt{sameSite=lax}, valable sept jours. Le jeton encode l'identifiant utilisateur, son
rôle et sa date d'expiration ; il est revérifié (signature et expiration) à chaque requête
via la fonction \texttt{getSession()}.

Le contrôle d'accès est appliqué à deux niveaux, conformément au principe de défense en
profondeur~: un middleware Next.js redirige les utilisateurs non autorisés au niveau du
routage, et chaque Server Action rappelle explicitement \texttt{requireRole([...])} en
première ligne avant tout traitement. Cette redondance garantit qu'aucune action
sensible ne peut être déclenchée par contournement de l'interface.

L'interface s'adapte au rôle connecté~: la barre latérale n'affiche que les sections
accessibles au rôle de l'utilisateur courant (par exemple, le Vigile ne dispose d'aucune
entrée de menu, son accès étant restreint à l'écran de vérification).

\capture{01-login}{0.6}{Écran de connexion}

\section{Gestion des clients}

Le module Clients permet de créer, consulter, modifier et supprimer des fiches client,
avec une distinction explicite entre trois types~: particulier, institution étatique et
entreprise. Les champs \texttt{raisonSociale} et \texttt{ice} deviennent obligatoires de
façon conditionnelle selon le type sélectionné, une règle de validation portée par un
schéma Zod avec logique de raffinement (\texttt{superRefine}).

La liste des clients propose une recherche par nom ou téléphone et une pagination côté
serveur. La suppression d'un client est bloquée métier si celui-ci possède encore des
véhicules enregistrés, afin de préserver l'intégrité référentielle.

\capture{02-clients-liste}{0.85}{Liste des clients avec recherche et pagination}
\capture{03-client-detail}{0.85}{Fiche client avec ses véhicules rattachés}

\section{Gestion des véhicules}

Le module Véhicules permet de rattacher un ou plusieurs véhicules à un client, avec les
caractéristiques usuelles (immatriculation, marque, modèle, année, couleur,
kilométrage). L'immatriculation est contrainte à l'unicité au niveau de la base de
données, et cette contrainte est explicitement vérifiée en amont dans le service métier
pour renvoyer un message d'erreur clair plutôt qu'une erreur technique.

La \textbf{recherche par immatriculation} a fait l'objet d'un soin particulier, identifiée
dès la spécification comme la fonction la plus utilisée au quotidien par l'équipe~: elle
est disponible directement sur la liste des véhicules, mais également sous la forme d'un
champ de recherche rapide intégré à la barre latérale, accessible depuis n'importe quelle
page de l'application et redirigeant vers les résultats correspondants.

La fiche détaillée d'un véhicule présente son historique complet des dossiers
d'intervention passés et en cours, avec un accès direct à la création d'un nouveau
dossier pour ce véhicule.

\capture{04-vehicules-liste}{0.85}{Liste des véhicules avec recherche par immatriculation}
\capture{05-vehicule-detail}{0.85}{Fiche véhicule avec historique des dossiers}

\section{Création et suivi des dossiers d'intervention}

Le module Dossiers constitue le cœur du système. Le formulaire de création propose deux
modes complémentaires, activables par bascule~: rattacher le dossier à un véhicule
\textbf{déjà enregistré}, ou créer \textbf{à la volée} un nouveau client et son véhicule
dans le même formulaire --- un cas fréquent en pratique lorsqu'un nouveau client se
présente directement au garage. La sélection du ou des types d'intervention est multiple
(cases à cocher), reflétant fidèlement la règle métier selon laquelle un dossier peut
cumuler plusieurs interventions.

\capture{08-dossier-nouveau-existant}{0.85}{Formulaire de création --- véhicule existant}
\capture{09-dossier-nouveau-client}{0.85}{Formulaire de création --- nouveau client et véhicule}

\subsection{Numérotation séquentielle sans collision}

Conformément à la règle métier n°7, chaque dossier reçoit un numéro séquentiel de la
forme \texttt{DOS-2026-XXXX}, calculé à partir du dernier numéro existant pour l'année en
cours. Pour garantir l'absence de collision en cas de créations concurrentes, l'ensemble
de l'opération --- résolution ou création du véhicule et du client, calcul du numéro,
création du dossier et écriture de l'entrée d'audit initiale --- est exécuté dans une
\textbf{transaction Prisma en niveau d'isolation \texttt{Serializable}}. En cas de conflit
d'écriture détecté par PostgreSQL, la transaction est automatiquement retentée (jusqu'à
cinq tentatives), garantissant à la fois l'absence de doublon de numéro et l'atomicité de
la création du client, du véhicule et du dossier associés.

\subsection{Fiche d'intervention}

Dès la création du dossier, une \textbf{fiche d'intervention} est généré à la demande au
format PDF (format A5), via \texttt{@react-pdf/renderer}, exposée par une route API
dédiée (\texttt{/api/dossiers/[id]/fiche}). Le document reprend le numéro de dossier,
l'immatriculation en gros caractères, la marque et le modèle, le client, le ou les types
d'intervention et le motif déclaré. Il s'agit d'un document purement informatif, sans QR
code ni valeur d'autorisation, destiné à être imprimé et placé physiquement dans le
véhicule.

\capture{06-dossiers-liste}{0.85}{Liste des dossiers avec filtres par statut et recherche}
\capture{07-dossier-detail}{0.85}{Fiche détaillée d'un dossier, avec accès à la fiche PDF}

% ================================================================
\chapter{État d'avancement et planning}
% ================================================================

Le développement est organisé sur quatre semaines, à raison d'une tâche métier par jour
ouvré, chacune donnant lieu à un commit atomique. Le tableau suivant synthétise le
planning global.

\begin{table}[h!]
\centering
\rowcolors{2}{km0gray}{white}
\begin{tabularx}{\linewidth}{@{}l>{\raggedright\arraybackslash}X@{}}
\toprule
\rowcolor{km0blue}
\textcolor{white}{\textbf{Semaine}} & \textcolor{white}{\textbf{Contenu}} \\
\midrule
\textbf{Semaine 1 --- Fondations} & Setup projet, authentification \& rôles, module clients,
module véhicules, création de dossier. \\
\textbf{Semaine 2 --- Cœur métier} & Vue Kanban de l'atelier, machine à états et audit,
génération du Bon de Sortie, écran de vérification, facturation. \\
\textbf{Semaine 3 --- Complétude et finition} & Tableau de bord, pièces \& détail dossier,
administration, polish UI/UX, corrections. \\
\textbf{Semaine 4 --- Tests et livraison} & Tests métier, tests de parcours de bout en bout,
données de démonstration, documentation, finalisation. \\
\bottomrule
\end{tabularx}
\caption{Planning de développement (détail complet dans \texttt{TASKS.md})}
\end{table}

\begin{table}[h!]
\centering
\rowcolors{2}{km0gray}{white}
\begin{tabularx}{\linewidth}{@{}>{\raggedright\arraybackslash}X l@{}}
\toprule
\rowcolor{km0blue}
\textcolor{white}{\textbf{Module}} & \textcolor{white}{\textbf{État}} \\
\midrule
Setup projet, schéma Prisma, données de démonstration & \etat{FAIT}{km0green} \\
Authentification par session et contrôle d'accès par rôle & \etat{FAIT}{km0green} \\
Gestion des clients & \etat{FAIT}{km0green} \\
Gestion des véhicules et recherche par immatriculation & \etat{FAIT}{km0green} \\
Création et listing des dossiers d'intervention & \etat{FAIT}{km0green} \\
Vue Kanban de l'atelier, machine à états, audit & \etat{À VENIR}{km0orange} \\
Génération du Bon de Sortie (PDF + QR) & \etat{À VENIR}{km0orange} \\
Écran de vérification de sortie & \etat{À VENIR}{km0orange} \\
Facturation et dérogation de paiement & \etat{À VENIR}{km0orange} \\
Tableau de bord de synthèse & \etat{À VENIR}{km0orange} \\
Administration et régénération de BS perdu & \etat{À VENIR}{km0orange} \\
Tests métier et données de démonstration finales & \etat{À VENIR}{km0orange} \\
\bottomrule
\end{tabularx}
\caption{État d'avancement à la date de rédaction (7 août 2026)}
\end{table}

\begin{avancement}{portée de ce rapport}
Ce rapport correspond à l'état d'avancement du projet à sa date de rédaction, soit la fin
de la première semaine de développement. Les chapitres de conception couvrant les modules
non encore réalisés (Kanban, Bon de Sortie, écran de vérification, facturation, tableau
de bord) seront complétés et, le cas échéant, corrigés dans les versions ultérieures de ce
document au fur et à mesure de leur implémentation effective, afin que le rapport final
reflète fidèlement le système livré.
\end{avancement}

% ================================================================
\chapter{Tests et démarche qualité}
% ================================================================

\section{Vérification statique}

Le mode strict de TypeScript est activé sur l'ensemble du projet, sans autorisation du
type \texttt{any}. La commande \texttt{tsc --noEmit} est exécutée avant chaque commit et
doit passer sans erreur ; c'est une condition bloquante du flux de travail décrit au
chapitre~1.

\section{Validation fonctionnelle manuelle}

Chaque module livré a été vérifié en conditions réelles~: lancement de l'application
(serveur de développement Next.js connecté à une instance PostgreSQL locale) et parcours
des scénarios nominaux et des cas limites directement dans le navigateur --- création,
modification, suppression, recherche, pagination, et rejet des cas invalides
(immatriculation en doublon, champs obligatoires manquants). Cette validation a permis de
détecter et corriger, avant livraison, des anomalies qui n'auraient pas été révélées par
la seule vérification de types.

\section{Stratégie de tests automatisés (planifiée)}

Conformément au planning (semaine 4), une suite de tests automatisés sera constituée avec
\textbf{Vitest}, couvrant en priorité la logique métier critique~: transitions de statut
autorisées, génération des numéros séquentiels, règles de génération et de vérification du
Bon de Sortie. Des scénarios de bout en bout viendront ensuite valider le parcours complet
d'un véhicule, de son entrée jusqu'à sa sortie.

% ================================================================
\chapter{Conclusion générale et perspectives}
% ================================================================

Le projet KM0 répond à un besoin concret observé sur le terrain~: l'absence d'outillage
numérique adapté à la gestion quotidienne d'un garage automobile traitant à la fois des
particuliers et des institutions étatiques. La démarche suivie --- analyse du contexte,
spécification formalisée des besoins, conception d'une architecture découplée et d'un
modèle de données couvrant l'intégralité du cycle de vie d'un dossier --- a permis de
poser des fondations solides, aujourd'hui concrétisées par un socle applicatif
opérationnel~: authentification et contrôle d'accès par rôle, gestion des clients et des
véhicules, et création de dossiers d'intervention avec génération automatique de la fiche
d'intervention.

Les prochaines étapes du développement porteront sur le cœur métier le plus sensible du
projet~: la vue Kanban de l'atelier et la machine à états des dossiers, puis la
génération et la vérification du Bon de Sortie, mécanisme central de sécurisation de la
sortie des véhicules. Viendront ensuite la facturation, le tableau de bord de synthèse et
les fonctionnalités d'administration, avant une phase dédiée aux tests automatisés et à la
préparation des données de démonstration.

Au-delà du périmètre du présent stage, plusieurs évolutions restent envisageables à plus
long terme~: un mode de fonctionnement en réseau local pour les sites sans connexion
Internet stable, une gestion de stock détaillée des pièces détachées, ou encore une
extension multi-garage. Ces évolutions n'ont pas été anticipées dans le code au-delà du
découplage naturel des services métier, conformément au principe retenu de ne pas
sur-architecturer une version dont le périmètre reste, à ce stade, celui d'un garage
unique.

% --------------------------------------------------------------
\chapter*{Bibliographie et webographie}
\addcontentsline{toc}{chapter}{Bibliographie et webographie}

\begin{itemize}[leftmargin=1cm]
\item Documentation officielle Next.js --- \url{https://nextjs.org/docs}
\item Documentation officielle Prisma --- \url{https://www.prisma.io/docs}
\item Documentation officielle Zod --- \url{https://zod.dev}
\item Documentation officielle Tailwind CSS --- \url{https://tailwindcss.com/docs}
\item Documentation MDN --- Web Crypto API --- \url{https://developer.mozilla.org/fr/docs/Web/API/Web_Crypto_API}
\item Documentation \texttt{@react-pdf/renderer} --- \url{https://react-pdf.org}
\item Documentation PostgreSQL 16 --- \url{https://www.postgresql.org/docs/16/}
\end{itemize}

% ================================================================
\chapter{Annexe --- Glossaire}
% ================================================================

\begin{table}[h!]
\centering
\rowcolors{2}{km0gray}{white}
\begin{tabularx}{\linewidth}{@{}l>{\raggedright\arraybackslash}X@{}}
\toprule
\rowcolor{km0blue}
\textcolor{white}{\textbf{Terme}} & \textcolor{white}{\textbf{Définition}} \\
\midrule
\textbf{Dossier} & Fiche représentant une intervention sur un véhicule, du dépôt à la sortie. \\
\textbf{Bon de Sortie (BS)} & Document autorisant la sortie physique d'un véhicule, muni d'un QR code à usage unique. \\
\textbf{Fiche d'intervention} & Document résumé généré à l'entrée du véhicule, placé physiquement dans celui-ci, informant le technicien du travail demandé. \\
\textbf{Dérogation} & Autorisation exceptionnelle accordée par un Admin pour sortir un véhicule non soldé. \\
\textbf{Vigile} & Agent de sécurité chargé de contrôler les sorties de véhicules à la guérite. \\
\textbf{Kanban} & Vue en colonnes représentant les dossiers selon leur statut d'avancement. \\
\textbf{Institution étatique} & Client de type organisme public (véhicules de fonction, administrations). \\
\textbf{Server Action} & Fonction serveur Next.js invocable directement depuis un composant React, sans passer par une route API explicite. \\
\bottomrule
\end{tabularx}
\caption{Glossaire des termes métier et techniques}
\end{table}

\end{document}
